\documentclass[10pt,twocolumn]{article}
\usepackage[margin=1.5cm]{geometry}
\usepackage{amsmath, amssymb}
\usepackage{titlesec}

\setlength{\parindent}{0pt}
\setlength{\columnsep}{0.8cm}
\setlength{\parskip}{3pt}

\titlespacing*{\section}{0pt}{10pt}{6pt}
\titlespacing*{\subsection}{0pt}{6pt}{4pt}

\begin{document}

\title{\vspace{-0.5cm}Burnside y Pólya -- Resumen}
\date{}
\maketitle
\vspace{-0.5cm}

\section*{Teoría}

Sirve para contar cuántos objetos realmente distintos existen cuando
algunos se consideran iguales bajo ciertas operaciones
(rotaciones, simetrías, intercambios, etc.).

Sea \(G\) un grupo que actúa sobre un conjunto.
Sea \(C(\pi)\) el número de ciclos de la permutación \(\pi\).

\subsection*{Lema de Burnside}

\[
|\text{órbitas}| = \frac{1}{|G|}\sum_{\pi\in G} I(\pi)
\]

\(I(\pi)\): configuraciones invariantes bajo \(\pi\).  
La respuesta es el promedio de configuraciones fijas.

\subsection*{Teorema de Pólya ($k$ colores)}

\[
|\text{órbitas}| = \frac{1}{|G|}\sum_{\pi\in G} k^{C(\pi)}
\]

Cada ciclo puede colorearse independientemente.

% -------------------------------------------------------------

\section*{CSES 2209 -- Necklaces}

Collares de \(n\) perlas con \(k\) colores,
equivalentes por rotación.

Ejemplo de igualdad:

Para \(n=4\),

\[
ABCD \sim BCDA \sim CDAB \sim DABC
\]

Todas son la misma rotación.

Grupo \(C_n\), \(|G| = n\).

\[
\boxed{
\frac{1}{n}\sum_{i=0}^{n-1} k^{\gcd(i,n)}
=
\frac{1}{n}\sum_{d\mid n}\phi\!\left(\frac{n}{d}\right)k^d
}
\]

% -------------------------------------------------------------

\section*{CSES 2210 -- Grids}

Grillas binarias \(n\times n\),
equivalentes por rotaciones.

Ejemplo de igualdad:

\[
\begin{matrix}
1 & 0 \\
0 & 1
\end{matrix}
\quad\sim\quad
\begin{matrix}
0 & 1 \\
1 & 0
\end{matrix}
\]

Son iguales bajo rotación de 90°.

Grupo \(C_4\), \(k=2\).

\[
a=\left(\frac{n}{2}\right)\left(\frac{n+1}{2}\right),
\qquad
o=n\bmod 2
\]

\[
\boxed{
\frac{1}{4}
\left(
2^{n^2}
+
2^{2a+o}
+
2\cdot 2^{a+o}
\right)
}
\]

% -------------------------------------------------------------

\section*{CF 1065E -- Transmutation}

Cadena de longitud \(n\) sobre alfabeto de tamaño \(k\).

Para cada \(b_i\): invertir prefijo y sufijo de longitud \(b_i\) e intercambiarlos.

Ejemplo simple:

Sea \(n=6\), \(b_1=2\).

\[
\underline{AB} \, C D \, \underline{EF}
\quad\longrightarrow\quad
\underline{FE} \, C D \, \underline{BA}
\]

Ambas cadenas son equivalentes.

\[
1\le b_1<\dots<b_m
\]

\[
d_1=b_1,\qquad d_i=b_i-b_{i-1}
\]

\textbf{Estructura.}

Solo afectan \(2b_m\) extremos.  
Centro libre:

\[
k^{\,n-2b_m}
\]

Cada capa \(d_i\) es independiente:

\[
\text{identidad: } k^{2d_i}
\qquad
\text{swap: } k^{d_i}
\]

Podemos usar o no cada \(b_i\) ⇒ grupo de tamaño \(2^m\).

\textbf{Burnside:}

\[
\boxed{
k^{\,n-2b_m}
\prod_{i=1}^{m}
\frac{k^{2d_i}+k^{d_i}}{2}
}
\]

\[
\boxed{
k^{\,n-b_m}
\cdot
\frac{1}{2^{m}}
\prod_{i=1}^{m}(1+k^{d_i})
}
\]

\end{document}

